\section{研究背景与任务目标}
\label{sec:introduction}

潮流计算用于确定给定网络拓扑、发电计划和负荷水平下各母线电压幅值、相角及线路功率分布，
是安全校核、稳定分析和优化调度的基础。

重载、弱网、长距离辐射供电或参数尺度差异明显时，潮流方程可能无实数解，
也可能虽有解但迭代无法进入吸引域，还可能因雅可比矩阵病态而出现数值"出轨"。
因此，本任务从物理、数学和计算三个层面分析不收敛问题。

\textbf{本项目的主要任务包括：}
\begin{enumerate}
  \item 解析 IEEE 测试数据并建立母线、支路和节点导纳矩阵模型；
  \item 实现并比较牛顿--拉夫逊法与快速解耦潮流法；
  \item 开展精度、初值、负荷增长和线路参数鲁棒性测试；
  \item 测试最优乘子法与非线性规划法在临界工况下的发散与残差停滞；
  \item 测试不同参数化方法对连续潮流法跟踪 $P$--$V$ 曲线及越过鼻点能力的影响。
\end{enumerate}

\textbf{病态不收敛问题的系统认识}
通过文献调研与理论学习，本报告将不收敛区分为三个层面：
\begin{enumerate}
  \item \textbf{物理层面——解是否存在：}
  负荷增长使系统逼近静态电压稳定极限；
  在鞍结分岔点处，稳定与不稳定平衡点重合，潮流雅可比矩阵趋于奇异。
  越过极限后，原控制方式下可能不存在对应的实数平衡解。
  \item \textbf{数学层面——能否找到解：}
  NR 等局部求根算法具有有限吸引域，即使方程存在解，
  不合适的初值仍可能导致振荡、越界或收敛到其他解支。
  \item \textbf{计算层面——迭代是否稳定：}
  雅可比矩阵条件数恶化会放大舍入误差，
  过大的修正步还会使迭代偏离有效区域，
  因此需要尺度化、稳定分解和步长控制。
\end{enumerate}

\textbf{病态不收敛问题的应对方法}
针对不同的不收敛成因，有不同的处理手段：
\begin{itemize}
  \item \textbf{物理层面——逼近物理解边界：}
  连续潮流法引入延拓参数扩展潮流方程，
  使迭代在雅可比奇异时仍能沿 $P$--$V$ 曲线追踪至鞍结分岔点，
  判断该工况下解是否真实存在。
  \item \textbf{数学层面——扩大吸引域：}
  非线性规划法将潮流方程转化为残差范数最小化问题，
  通过全局寻优降低对初值质量的依赖。
  \item \textbf{计算层面——改善数值稳定性：}
  标幺化消除量级差异（尺度化），
  带主元选取的 LU 分解避免小主元除零（稳定分解），
  最优乘子法沿 NR 方向一维寻优限制修正幅度（步长控制）；
  此外，迭代顺序与数值截断策略也对稳定性有一定影响。
\end{itemize}
